\documentclass[11pt]{article}
\usepackage[top=2.3cm, bottom=2.3cm, left=2.1cm, right=2.1cm]{geometry}
\usepackage{amsmath,amssymb}
\usepackage{booktabs}
\usepackage[T1]{fontenc}
\usepackage{lmodern}
\usepackage{hyperref}
\hypersetup{colorlinks=true, linkcolor=blue, urlcolor=blue, citecolor=blue}

\title{Proposition 2 for the Reference Spec}
\author{}
\date{\today}

\begin{document}

\maketitle

This note answers two questions for the reference firm regression
\begin{quote}
\texttt{firm\_\_main\_\_levels\_\_extensive\_\_coalition\_\_cycle\_specific\_\_unweighted\_\_all\_firms\_\_pooled\_count}
\end{quote}

\begin{enumerate}
  \item How can we show Proposition 2 really holds in the data?
  \item Why does the simple regression on cell means still give different coefficients?
\end{enumerate}

All numbers below come from:
\begin{itemize}
  \item \path{BNDES/output/diagnostics/prop2_reference_spec_comparison.csv}
  \item \path{BNDES/output/diagnostics/prop2_reference_spec_diagnostics.csv}
  \item \path{BNDES/output/diagnostics/prop2_reference_spec_cell_fe.qs2}
\end{itemize}

\section*{Short Answer}

Proposition 2 \textbf{does hold} if we start from the firm regression itself, keep the exact estimation sample that the firm regression used, save the fitted fixed effects, and aggregate the \textbf{firm sufficient statistics}.

Proposition 2 \textbf{does not hold} if we replace that object with a simpler new regression on cell means. That simpler regression loses important within-cell variation.

So there are really two different exercises:
\begin{itemize}
  \item \textbf{Exact aggregation test}: this is the right test of Proposition 2. It now works.
  \item \textbf{Simple cell-mean regression}: this is only an approximation. On real data it still differs.
\end{itemize}

\section{The Two Samples: Group A and Group B}

Define
\[
A = \{i: y_i \text{ and } x_i \text{ are observed}\},
\]
\[
B = \{i \in A: i \text{ is actually used by the firm regression}\}.
\]

Here:
\begin{itemize}
  \item $|A| = 44{,}181{,}405$
  \item $|B| = 43{,}184{,}682$
  \item $|A \setminus B| = 996{,}723$
\end{itemize}

So the firm regression does \textbf{not} use all rows in $A$. It uses only $B$.

This matters because the aggregate version must be built from $B$, not from $A$. If we aggregate $A$ but the firm regression used $B$, then the two sides are not even using the same data.

The new script now fixes this directly:
\begin{itemize}
  \item fully kept firms: $5{,}464{,}381$
  \item partially dropped firms: $121$
  \item fully dropped firms: $996{,}453$
  \item fully dropped singleton firms: $996{,}450$
\end{itemize}

In simple language: almost all dropped firms disappear completely from the firm regression, so they should also disappear from the aggregated version.

\section{The Exact Firm Equation}

On the used sample $B$, the firm regression is
\[
y_i = \beta x_i + FE_i + u_i,
\]
where:
\begin{itemize}
  \item $y_i = \mathbf{1}(\mathrm{BNDES}_{fmt} > 0)$
  \item $x_i$ is one of \texttt{FA\_mayor\_coalition}, \texttt{FA\_gov\_coalition}, \texttt{FA\_pres\_coalition}
  \item $FE_i$ is the fitted sum of the firm fixed effect and the municipality-year fixed effect
\end{itemize}

After estimating the firm regression, define
\[
y_i^* = y_i - \widehat{FE}_i.
\]

Then, with one regressor, the firm coefficient is
\[
\hat\beta_{\text{firm}} =
\frac{\sum_{i \in B} x_i y_i^*}{\sum_{i \in B} x_i^2}.
\]

This is the key formula.

\section{How to Show Proposition 2 Holds Exactly}

Partition the used sample $B$ into cells $c = (j,m,t)$. For each cell define
\[
S_{xy,c} = \sum_{i \in B_c} x_i y_i^*,
\qquad
S_{xx,c} = \sum_{i \in B_c} x_i^2.
\]

Then
\[
\hat\beta_{\text{firm}} =
\frac{\sum_c S_{xy,c}}{\sum_c S_{xx,c}}.
\]

This is just regrouping the same firm sums by cell. Nothing is approximated.

This is the exact way to show Proposition 2 holds in code:
\begin{enumerate}
  \item Run the firm regression.
  \item Save the exact used sample $B$.
  \item Save the fitted fixed effects $\widehat{FE}_i$.
  \item Compute $y_i^* = y_i - \widehat{FE}_i$.
  \item Aggregate $S_{xy,c}$ and $S_{xx,c}$ by cell.
  \item Rebuild $\hat\beta$ from those cell sums.
\end{enumerate}

That is exactly what the new script does.

\subsection*{Result}

\begin{center}
\begin{tabular}{lrrr}
\toprule
Tier & Firm beta & Exact beta from cell sums & Exact gap \\
\midrule
Mayor & \texttt{0.002176} & \texttt{0.002176} & \texttt{-2.20e-11} \\
Governor & \texttt{0.003404} & \texttt{0.003404} & \texttt{-6.56e-12} \\
President & \texttt{0.000378} & \texttt{0.000378} & \texttt{1.86e-12} \\
\bottomrule
\end{tabular}
\end{center}

The largest cell-level identity error is only \texttt{1.11e-16}.

So this part is now clear:
\begin{quote}
If we carry the exact firm regression solution to the cell level, Proposition 2 holds for the reference spec.
\end{quote}

\section{Why the Simple Regression on Cell Means Still Fails}

Now define cell means on the used sample $B$:
\[
\bar x_c = \frac{1}{N_c}\sum_{i \in B_c} x_i,
\qquad
\bar y_c^* = \frac{1}{N_c}\sum_{i \in B_c} y_i^*.
\]

The simple aggregated regression uses
\[
\hat\beta_{\text{mean}} =
\frac{\sum_c N_c \bar x_c \bar y_c^*}{\sum_c N_c \bar x_c^2}.
\]

This looks similar, but it is \textbf{not} the same object as the exact firm formula.

The reason is:
\[
\sum_{i \in B_c} x_i^2
=
N_c \bar x_c^2
+
\sum_{i \in B_c} (x_i - \bar x_c)^2,
\]

and
\[
\sum_{i \in B_c} x_i y_i^*
=
N_c \bar x_c \bar y_c^*
+
\sum_{i \in B_c} (x_i - \bar x_c)(y_i^* - \bar y_c^*).
\]

So the simple mean regression drops two within-cell terms:
\begin{itemize}
  \item within-cell variance of $x$
  \item within-cell covariance between $x$ and $y^*$
\end{itemize}

If those within-cell terms are small, the mean regression is close. If those terms are large, the mean regression is different.

In this dataset they are large.

\subsection*{Direct Evidence}

The share of the exact denominator that comes from the dropped within-cell part is:
\begin{itemize}
  \item Mayor: \texttt{91.9\%}
  \item Governor: \texttt{91.9\%}
  \item President: \texttt{93.7\%}
\end{itemize}

This means most of the identifying variation is \textbf{inside} cells, not between cell means.

That is why the simple mean regression still differs:

\begin{center}
\begin{tabular}{lrrr}
\toprule
Tier & Firm beta & Mean-regression beta & Gap \\
\midrule
Mayor & \texttt{0.002176} & \texttt{0.001290} & \texttt{-0.000886} \\
Governor & \texttt{0.003404} & \texttt{0.004397} & \texttt{0.000993} \\
President & \texttt{0.000378} & \texttt{-0.006601} & \texttt{-0.006979} \\
\bottomrule
\end{tabular}
\end{center}

So the mean regression is not failing because Proposition 2 is false. It is failing because cell means are not enough to summarize the firm-level problem.

\section{C1 to C6 in Simple Language}

Below is the clean way to read the conditions from the plan.

\subsection*{C1. Correct weighting}

Plain language: a cell with 100 firms should count 100 times more than a cell with 1 firm.

In the mean regression this means using weight \texttt{N\_c}.

Status here: \textbf{enforced}.

\subsection*{C2. Same sample}

Plain language: the aggregated version must use the exact same firm rows as the firm regression.

Here that means: aggregate $B$, not $A$.

Status here: \textbf{enforced}.

\subsection*{C3. Firm immobility / FE nesting}

Plain language: if the same firm appears in more than one cell, then one simple cell fixed effect cannot play the role of that firm's fixed effect.

Status here:
\begin{itemize}
  \item \textbf{not needed for the exact aggregation test}, because we keep the firm FE from the firm regression;
  \item \textbf{still a problem for a separate simple cell FE regression}.
\end{itemize}

\subsection*{C4. Correct aggregated FE object}

Plain language: the aggregated nuisance term must be the one implied by the firm regression, not a guessed simpler fixed effect.

Status here: \textbf{enforced in the exact test}, because we use the saved fitted FE from the firm regression and average it to the cell level.

\subsection*{C5. Fixed cell composition}

Plain language: if the firms inside a cell change over time, then the average firm fixed effect of that cell also changes over time.

Status here:
\begin{itemize}
  \item \textbf{not needed for the exact aggregation test}, because we work cell-year by cell-year using the saved firm solution;
  \item \textbf{still a problem for a separate simple cell FE regression}.
\end{itemize}

\subsection*{C6. No within-cell regressor heterogeneity after FE removal}

Plain language: once fixed effects are removed, firms inside the same cell should all have the same $x$ if cell means are supposed to be enough.

Status here: \textbf{fails on real data}, and this is the main reason the simple mean regression is still different.

\section{What Should Be Shown in the Paper or Presentation}

The clearest presentation is:
\begin{enumerate}
  \item Show the firm regression.
  \item Define $A$ and $B$, and state that the aggregation uses $B$.
  \item Show the exact identity
  \[
  \hat\beta = \frac{\sum_c S_{xy,c}}{\sum_c S_{xx,c}}.
  \]
  \item Report that the exact cell-sum coefficient matches the firm coefficient up to machine precision.
  \item Then show the mean-regression coefficient separately and explain that it differs because the dropped within-cell terms are large.
\end{enumerate}

That tells the full story very simply:
\begin{quote}
Proposition 2 holds exactly when we aggregate the firm regression itself. A new regression on cell means is only an approximation, and on real data that approximation is not exact because within-cell variation is large.
\end{quote}

\end{document}
